% !PRINTED 105
\weberpage{105}
%<<CONT

\begin{equation}
 |f(x)|>C.
\tag{3}
\end{equation}
Hieran schliesst sich der Satz

\medskip
\noindent V. \gesperrt{$f(x)$ ist eine stetige Function von $x$.}
\medskip

Damit soll folgendes gesagt sein: Ist der absolute Werth von $x$ kleiner als eine beliebige endliche Grösse $R$, so kann man, wie klein auch die positive Grösse $\omega$ angenommen wird, eine positive Grösse $\varepsilon$ derart finden, dass
\begin{equation}
 |f(x+h)-f(x)|<\omega
\tag{4}
\end{equation}
wird, so lange der absolute Werth $\varrho$ von $h$ kleiner als $\varepsilon$ bleibt, und zwar so, dass $\varepsilon$ nur von der Wahl von $\omega$, nicht aber von $x$ abhängig ist.

Nach §. 13, (2) ist
\begin{equation}
\begin{split}
 f(x+h)-f(x)&=h f'(x)+\frac{h^2}{1\cdot2}f''(x)+\cdots+\frac{h^n}{\Pi(n)}f^{(n)}(x)\\
 &=h\left\{f'(x)+\frac{h}{2}f''(x)+\cdots+a_0 h^{n-1}\right\}.
\end{split}
\tag{5}
\end{equation}
Nehmen wir nun den absoluten Werth von $h$ kleiner als $1$ an, so ist der absolute Werth der in der Klammer stehenden Grösse
\begin{equation}
 f'(x)+\frac{h}{2}f''(x)+\cdots+a_0 h^{n-1}
\tag{6}
\end{equation}
kleiner als eine bestimmte von Null verschiedene Zahl $k$, die man erhält, wenn man in (6) $h$ durch $1$, die Coëfficienten $a_0,a_1,\ldots,a_{n-1}$ durch ihre absoluten Werthe und $x$ durch $R$ ersetzt, weil dadurch jedes einzelne Glied der Summe (6) durch seinen absoluten Werth oder durch eine grössere Zahl ersetzt wird. Wählt man $\varepsilon$ so klein, dass
\[
 \varepsilon k<\omega
\]
ist, so ist wegen (5) die Forderung (4) erfüllt, so lange
\[
 \varrho<\varepsilon
\]
bleibt.

Wir können dem Satze von der Stetigkeit der Function $f(x)$ noch einen etwas anderen Ausdruck geben, der uns für die Folge wichtig ist. Nach dem Satze, dass der absolute Werth einer Summe von zwei Gliedern der Grösse nach zwischen der Summe und der Differenz der absoluten Werthe der Glieder liegt (vgl. Einleitung, S. 19), erhalten wir aus (5), wenn wir
%CONT>>
